\section{Introduction}

Modern machine learning systems depend on data infrastructure that can store
features, update them from new events, and recover from failures. The goal of
MLStore-Lite is to build a small educational prototype of such a system. The
project is inspired by the distributed data-system concepts from \textit{Designing
Data-Intensive Applications}, but keeps the implementation local and readable.

The project does not attempt to compete with production systems such as RocksDB,
Cassandra, Kafka, Spark, or Flink. Instead, it implements compact versions of
their central ideas. This makes the internal architecture easier to inspect:
write-ahead logging, sorted-string tables, replication, partitioning, batch
processing, stream processing, and observability can all be studied in one code
base.

The final system can run a small feature workflow end to end. Raw historical
events are converted into batch features, new events are appended to a stream
log, stream consumers update windowed features, and all feature values are
stored in a sharded replicated key-value store. The last layer adds local
evaluation and observability so the system can be measured and explained.

\section{System Overview}

MLStore-Lite is organized as a layered system. Each layer has a narrow
responsibility and builds on the layer below it:

\begin{verbatim}
Storage engine
  -> Replication
  -> Sharding / partitioning
  -> Batch processing
  -> Stream processing
  -> Integration
  -> Evaluation and observability
\end{verbatim}

This structure is useful because it separates concerns. The storage engine only
needs to know how one node persists key-value data. The replication layer only
needs to know how several nodes keep copies of the same data. The sharding layer
only decides which replicated group owns a key. Batch and stream processing then
use the sharded store without needing to understand the lower-level storage
details.

\subsection{Storage Engine}

The storage engine answers the question: how does one node store data durably?
It consists of a write-ahead log, an in-memory table, immutable SSTable files,
and compaction.

When a key is written, the value is first appended to the write-ahead log. This
means the operation can be recovered even if the process crashes before the
in-memory data is flushed to disk. The in-memory table keeps recent writes fast.
When it grows large enough, it is flushed into an SSTable file. SSTables are
immutable sorted files, which makes reads and merges more predictable.
Compaction later combines SSTables and removes overwritten values.

This is a simplified log-structured storage engine. The same broad idea appears
in systems such as LevelDB and RocksDB, although production engines add many
optimizations such as bloom filters, compression, concurrency control, and more
advanced compaction strategies.

\subsection{Replication}

The replication layer answers the question: how can the system keep multiple
copies of the same data? MLStore-Lite uses a leader-follower model. Each
replicated group contains one leader and two followers, giving a replication
factor of three.

Writes go to the leader. The leader stores the write locally and forwards it to
followers. Reads are served through the current leader in the implemented
interface. The project also supports manual failover, where a follower can be
promoted to leader after a failure experiment.

Replication improves fault tolerance because one physical copy is no longer the
only place where data exists. In this prototype the replicas are local objects
and directories, not separate networked machines. This keeps the focus on the
data-system concept rather than on networking.

\subsection{Sharding and Partitioning}

Replication creates copies of the same data. Sharding divides different keys
across different replicated groups. This answers the question: which shard owns
this key?

MLStore-Lite uses consistent hashing for routing. A key is hashed onto a ring,
and the next shard on the ring owns that key. This means keys can be distributed
without manually assigning ranges such as ``users 1--1000 go to shard A''. When
a shard is added, only part of the key space needs to move.

In the implemented topology, each shard is itself a replicated group with
replication factor three. Therefore a shard is not just one file or one Python
object. It is a logical partition of the key space, backed by multiple replica
nodes.

\subsection{Batch Processing}

The batch layer answers the question: how can raw historical records be turned
into useful derived data? In MLStore-Lite, the derived data are machine-learning
style features.

The batch engine follows a small MapReduce pattern:

\begin{verbatim}
map -> shuffle -> reduce
\end{verbatim}

The map step reads one input event and emits intermediate key-value pairs. The
shuffle step groups intermediate values by feature key. The reduce step combines
the grouped values into final feature values. For example, several click events
for the same user can become a single click-count feature.

The implemented batch feature job computes per-user event counts, click counts,
purchase counts, and total purchase amount. The results are written back into
the sharded replicated store using keys such as:

\begin{verbatim}
feature:user:42:click_count
\end{verbatim}

This mirrors the role of batch systems such as MapReduce and Spark, but at a
small local scale.
